\documentclass[11pt]{article}
\usepackage[a4paper,margin=26mm]{geometry}
\usepackage{amsmath,amssymb,amsthm,mathrsfs,hyperref}
\setlength{\emergencystretch}{3em}
\newtheorem{theorem}{Theorem}
\newtheorem{lemma}[theorem]{Lemma}
\newtheorem{proposition}[theorem]{Proposition}
\newcommand{\PSH}{\operatorname{PSH}}
\newcommand{\vol}{\operatorname{vol}}
\newcommand{\EE}{\mathcal E^1_\theta}
\newcommand{\OO}{\mathcal O}
\title{Actual outer section energy and its full-volume derivative}
\author{Unreviewed research, continuation 207}
\date{6 October 2026}
\begin{document}
\maketitle
\begin{abstract}
On every smooth projective variety over an algebraically closed
trivially valued characteristic-zero field, and in every actual big
rational bundle, we construct the outer section energy on the literal
single-net psh cone. Finite divisorial supporting planes give its
concavity, weak upper semicontinuity and decreasing-net continuity.
The independently checked uniform finite-model variation constructs
a target-independent positive Radon derivative of full volume on
the entire actual finite-energy domain. A two-sided touching argument
identifies the derivative at a model envelope. It is continuous along
arbitrary weak nets with converging energies and misses every proper
algebraic subset. The only non-Archimedean inputs are the already
checked actual scalar cone and the finite-model theorems, whose exact
scopes are stated below. No unresolved whole-energy or whole-operator
signature is granted. Unbounded physical weighted moving products
and the complete original equation are not conclusions of this paper.
\end{abstract}

\section{Domain and precise mathematical inputs}
Let $X$ be a smooth integral projective $n$-fold, $n\geq1$, over
an algebraically closed field $k$ of characteristic zero with the
trivial absolute value. Let $L\in\operatorname{Pic}(X)\otimes\mathbb Q$
be an actual big rational bundle. Set
$K=X^{\mathrm{an}}$, $\theta=c_1(L)$ and $V=\vol_X(L)>0$.
Retain all supported semivaluations on the compact Hausdorff carrier,
all directed nets, and actual bundle identities when taking sections.
Let $o$ be the generic trivial valuation. Canonical logarithmic norms
of regular sections are nonpositive and vanish at $o$.

An FS scalar is $m^{-1}\max_j(\log|s_j|+\lambda_j)$ on an
actual power of a rational bundle. Common zeros are allowed.
The genuine cone consists of all-point decreasing limits of one
FS net whose actual bundle classes tend to $\theta$, with finite
generic endpoint. Its weak topology is coordinate convergence at
rational divisorial valuations and $o$. We use the following
previously independently checked results in their exact scopes.
\begin{enumerate}
\item The actual scalar result of \cite[Sections 2--4]{Scalar}:
this cone is convex, finite-max and decreasing closed, has genuine
all-point envelopes $P_\theta(g)$ for upper bounded usc obstacles
with a genuine minorant, and has compact normalized weak topology.
Divisorial coordinates are finite and determine all-point order.
The genuine reference $R=P_\theta(0)$ has $R(o)=0$ and
every normalized genuine $u$ satisfies $u\leq R\leq0$.
Rational continuous model functions form a uniformly dense lattice
in $C(K)$. This result was reconstructed from the complete actual
ample grant G0 and relative generation grant G1; it grants no energy
or big-class equation.
\item \cite[Section 4]{Finite}, only its unconditional finite-model
conclusions: full section-lattice energy, its positive finite
divisorial derivative, refinement compatibility, positive deficient
derivatives and the uniform signed Taylor estimate displayed in
Section~\ref{variation}. These hold for bounded rational model
obstacles in the actual big $L$, not for an unbounded output.
No conditional S--M clause or inverse equation from that paper is
an input here.
\end{enumerate}
These two whole proof bodies were independently reviewed together
in round 161. Their accepted mathematical scopes, rather than their
historical manuscript status labels, are used. The algebraic
ingredients of the finite result are complete section-volume
asymptotics, volume differentiability and positive-product order on
smooth projective determinations, and normal birational section
descent. Their hypotheses remain the ones just stated. In particular,
there is no finite-generation or field-countability assumption.

\section{Full section traces and continuous energy}
We spell out the finite input to fix the energy being extended.
For a smooth equivariant model $Y\to X\times\mathbb P^1$,
product away from zero, write the full fiber as $F=\sum_E b_EE$.
A rational vertical Cartier divisor $D$ represents a continuous
model function $f_D$ with $f_D(x_E)=\operatorname{coeff}_E(D)/b_E$.
Choose $aF\leq D\leq bF$. In cleared degrees, the full relative
section lattice on $R_m=H^0(X,mL)$ is
\[
 t^{-ma}R_m[t]\subset\Lambda_m
 =\bigoplus_jt^{-w_{mj}}k[t]s_{mj}\subset t^{-mb}R_m[t].
                                                        \tag{1}
\]
Normal proper birational descent retains the full $R_m$.
Gluing to its product lattice at infinity gives
\[
 h^0(Y,m(\rho_X^*L+D+cF))
       =\sum_j\max(w_{mj}+mc+1,0).
                                                        \tag{2}
\]
For sufficiently large fixed rational $c$, section asymptotics imply
\[
 \frac{\vol_Y(\rho_X^*L+D+cF)}{n+1}
       =cV+\mathscr E(f_D),\qquad
 \mathscr E(f_D)=\lim_m\frac{n!}{m^{n+1}}\sum_j w_{mj}.
                                                        \tag{3}
\]
The limit is over all sufficiently large cleared admissible degrees.
The $+1$ terms are rank-order errors $O(m^n)$.
The finite result proves refinement invariance, monotonicity,
the constant law $\mathscr E(f+c)=\mathscr E(f)+cV$, and
$V$-Lipschitz continuity. Volume-root concavity applied to (3),
followed by expansion as the common time twist tends to infinity,
proves ordinary concavity. Thus $\mathscr E$ has a unique continuous
extension to $C(K)$ with these properties.

For clarity this extension is an actual full-section trace.
For $f\in C(K)$ and $0\ne s\in R_m$ define
\[
 W_m(f;s)=\inf_K(mf-\log|s|),\qquad
 e_m(f)=\frac{n!}{m^{n+1}}
                \sum_j\lfloor W_m(f;s_{mj})\rfloor,       \tag{4}
\]
where the basis is adapted to the indicated filtration.
At a section pole the bracket is $+\infty$; at $o$ it is finite.
Since $\log|s|\leq0$, every weight is finite.
The ultrametric inequality gives
$W(s+t)\geq\min(W(s),W(t))$ and invariance under nonzero scalar
multiplication. Vectors with distinct weights are linearly independent:
in any nontrivial linear relation the term of least weight is unique,
and cannot cancel. There are consequently at most $\dim R_m$
distinct weights. Their superlevel spaces give a finite filtration
with attained endpoints and an adapted basis. Its trace is intrinsic.
For a model obstacle, section extension identifies its integer
weights with (1). Comparison of every filtration space gives
\[
 |e_m(f)-e_m(g)|
 \leq\frac{n!\dim R_m}{m^n}
                   (\|f-g\|_\infty+1/m).               \tag{5}
\]
Uniform model density, (3), and
$n!\dim R_m/m^n\to V$ show $e_m(f)\to\mathscr E(f)$ for every
continuous $f$. Thus no abstract extension has replaced the full
section trace. For $f=0$, all weights are zero, so $\mathscr E(0)=0$.

The finite derivative is
\[
 \mu_D=\sum_E b_E\langle(\rho_X^*L+D+cF)^n\rangle\cdot E
                 \,\delta_{x_E}\geq0,
 \qquad \mu_D(K)=V.                                    \tag{6}
\]
Every coefficient is independent of sufficiently large $c$.
The effective-divisor volume derivative makes it nonnegative;
the fiber derivative gives the raw mass $V$.
For an exceptional prime on a higher smooth determination,
$p_*\OO(rG)=\OO$ for integral $r\geq0$, so its derivative
coefficient at the pulled big class is zero.
This and pullback compatibility preserve (6) under refinement.
For each model direction $h$,
$D\mathscr E_{f_D}(h)=\mu_D(h)$.
Concavity and uniform model approximation therefore give the global
supporting plane
\[
 A_D(f)=\mathscr E(f_D)+\mu_D(f-f_D),\qquad
 \mathscr E(f)\leq A_D(f)\quad(f\in C(K)).              \tag{7}
\]
If $f_D$ uniformly approximates $f$, their difference in (7) is
at most $2V\|f_D-f\|_\infty$. Consequently
\[
                    \mathscr E(f)=\inf_D A_D(f).        \tag{8}
\]
The plane index ranges over every rational model obstacle and is
independent of $f$.

\section{Actual outer energy, with no whole-energy grant}
Define on the actual genuine cone
\[
 E(u)=\inf_{f\in C(K),\ f\geq u}\mathscr E(f),\qquad
 \EE=\{u\in\PSH(\theta):E(u)>-\infty\}.                \tag{9}
\]
The energy is monotone and $E(u+c)=E(u)+cV$.

\begin{lemma}[Envelope identity and decreasing continuity]\label{outer}
For every $f\in C(K)$, $E(P_\theta(f))=\mathscr E(f)$.
For every decreasing genuine net $u_i\downarrow u$ with finite
generic endpoint, $E(u_i)\downarrow E(u)$, including $E(u)=-\infty$.
\end{lemma}
\begin{proof}
The FS scalar $m^{-1}(\log|s|+\lambda)$ is genuine on the actual $L$.
Whenever it lies below $f$, it lies below $P_\theta(f)$.
Thus, if continuous $g\geq P_\theta(f)$, every feasible weight
$\lambda$ for $f$ is feasible for $g$ in (4). All filtration
spaces for $f$ are contained in the corresponding ones for $g$,
so $\mathscr E(g)\geq\mathscr E(f)$ by (4).
Since $f$ itself majorizes its envelope, (9) gives equality.
This proof uses feasible inequalities, with no subtraction of two
infinite values at a common pole.

Now let continuous $f\geq u$ and $\epsilon>0$.
The closed sets $\{u_i\geq f+\epsilon\}$ are decreasing and
have empty intersection. Compactness of $K$ makes one tail empty.
On that tail $E(u_i)\leq\mathscr E(f)+\epsilon V$.
Monotonicity gives $\inf_iE(u_i)\geq E(u)$.
Taking the infimum over $f$, and then $\epsilon\downarrow0$,
proves the reverse inequality, also when this infimum is $-\infty$.
\end{proof}

In particular $E(R)=0$. Normalize $u$ at $o$ and put
$B(u)=-E(u)\geq0$. The genuine floors $\max(u,R-N)$ have cost
at most $NV$ and decrease all-point to $u$, including at common
poles. Lemma~\ref{outer} identifies $B(u)$ with their limiting
cost; bounded relative difference is not imposed on $\EE$.

Every upper bounded usc function on $K$ is the infimum of its
continuous majorants, directed by finite minima. Rational model
density gives model majorants with the same infimum. For a normalized
genuine $u\leq0$, taking their minima with zero gives nonpositive
model majorants $f_i\downarrow u$ with $f_i(o)=0$.
Here the model lattice and rational positive approximation margins
ensure that each approximant remains a majorant.
Their genuine envelopes decrease to $u$: the limit is genuine by
the scalar input, dominates $u$, and is bounded above by
$\inf_i f_i=u$. Lemma~\ref{outer} gives
$\mathscr E(f_i)\downarrow E(u)$.

\begin{theorem}[Finite planes on genuine and usc domains]\label{planes}
For every genuine $u$,
\[
                         E(u)=\inf_D A_D(u).            \tag{10}
\]
It is concave and weakly upper semicontinuous. The normalized
sublevels $B(u)\leq B<\infty$ are compact sets of actual functions
in $\EE$. If $g$ is upper bounded usc and finite at each divisorial
point, its outer continuous energy is likewise $\inf_D A_D(g)$.
\end{theorem}
\begin{proof}
Take continuous majorants $f_i\downarrow u$. They are cofinal for
the infimum in (9): for any other continuous majorant $f$ the
compact strict-margin argument gives eventually $f_i<f+\epsilon$.
For fixed $D$, its measure has finite divisorial support, where
$u$ is finite, so $\inf_i A_D(f_i)=A_D(u)$ by an ordinary finite
sum. Interchange the two infima in (8) to obtain (10).
The same proof applies to $g$. Each plane is affine and continuous
in finitely many weak coordinates, proving concavity and weak usc.
The energy superlevel is closed in the scalar normalized compact
cone. Every point on it has finite energy, proving compactness on
the actual domain. No measurable-net monotone convergence is used.
\end{proof}

For $u\in\EE$ and $h\in C(K)$, the envelope of $u+th$ exists:
$u-|t|\|h\|_\infty$ is a genuine minorant. Moreover
\[
\begin{split}
 H_{t,h}(u)&:=E(P_\theta(u+th))
       =\inf_D\{A_D(u)+t\mu_D(h)\},\qquad H_{0,h}(u)=E(u).
\end{split}                                             \tag{11}
\]
Indeed $P_\theta(f_i+th)\downarrow P_\theta(u+th)$ by the
same genuine closure and envelope order argument, and decreasing
energy continuity identifies its energy with the outer energy of
$u+th$. Theorem~\ref{planes} gives (11). Hence $H_{t,h}$ is weakly
usc at every fixed $t$, on the entire actual $\EE$.

\section{Uniform finite variation and the original derivative}\label{variation}
The unconditional finite Taylor theorem of \cite[Section 4]{Finite} states that
for fixed rational model $h$, $B\geq0$ and $c>B/V$ there are
$K_{B,c,h}<\infty$ and $\tau_{B,c,h}>0$, independent of the
nonpositive model obstacle $f_D$ and its determination, such that
\[
 |\mathscr E(f_D+th)-\mathscr E(f_D)-t\mu_D(h)|
   \leq K_{B,c,h}t^2+|t|a_{B,h}/c                       \tag{12}
\]
if $\mathscr E(f_D)\geq-B$ and $|t|<\tau_{B,c,h}$, where
$a_{B,h}=\|h\|_\infty(2B+2V\|h\|_\infty)$.
Its signed bigness margin is imposed before integration:
for $H_0=\|h\|_\infty>0$ take additionally
$|t|<(c-B/V)/(2H_0)$.
At parameter $s$ the nonpositive obstacle is
$f_D+sh-|s|H_0$, its cost is at most $B+2V|s|H_0$,
and its deficient time parameter is $c+|s|H_0$.
The margin makes this latter parameter larger than its cost divided
by $V$ for both signs. The proof uses two fixed nef tests for the
signed Cartier direction, and retains the positive mass tail
$\mu_D-\mu_{D,c}$, whose mass is at most $B/c$.
It does not assume an unbounded non-Archimedean Taylor theorem.

\begin{theorem}[Entire actual derivative]\label{derivative}
There is a uniquely constructed target-independent positive Radon
measure $M(u)$ on the original $K$ for every $u\in\EE$, of raw
mass $V$, such that for every absolute original $h\in C(K)$
\[
 \left.\frac{d}{dt}\right|_{0}E(P_\theta(u+th))=M(u)(h).
                                                        \tag{13}
\]
Both derivative signs hold. Adding a constant to $u$ preserves $M(u)$.
The uniform estimate (12) holds with $E$, $H_{t,h}$ and $M(u)$
on every normalized genuine cost sublevel.
\end{theorem}
\begin{proof}
Normalize $u$, take the nonpositive model majorants above, and put
$B=B(u)$. Their energies are at least $-B$.
The positive measures $\mu_{f_i}$ have raw mass $V$, so possess
weak clusters on compact $K$. Along a subnet with cluster $\lambda$,
decreasing energy continuity in (11) and (12) give
\[
 |H_{t,h}(u)-E(u)-t\lambda(h)|
       \leq K_{B,c,h}t^2+|t|a_{B,h}/c.                  \tag{14}
\]
Hold $c$ fixed before sending $t\to0$ from either sign.
The limsup and liminf of each difference quotient are within
$a_{B,h}/c$ of $\lambda(h)$. Then let $c\to\infty$.
Thus both quotients converge to $\lambda(h)$.
Every cluster has the same model-test integrals; uniform model
density and their common raw mass imply equality of the measures.
Compactness therefore gives whole-net convergence to a unique
$M(u)$ and the uniform genuine estimate (14).
The construction uses the actual energy, not an inverse target or
an abstract completion.

If $h\in C(K)$, approximate it uniformly by rational model $h_j$.
Envelope order, (9) and the raw constant law give
\[
 |H_{t,h}(u)-H_{t,h_j}(u)|
       \leq |t|V\|h-h_j\|_\infty.
\]
The same bound without $|t|$ holds for measure integrals.
First differentiate at fixed $j$, then $j\to\infty$, proving
(13). The constant law and
$P_\theta(u+b+th)=P_\theta(u+th)+b$ prove constant invariance.
\end{proof}

\begin{proposition}[Model-envelope compatibility]\label{touch}
For every rational model obstacle $f_D$, the derivative of its genuine
envelope is the finite measure:
\[
                         M(P_\theta(f_D))=\mu_D.
\]
\end{proposition}
\begin{proof}
Put $u=P_\theta(f_D)$. It belongs to $\EE$ by Lemma~\ref{outer}.
Since $u\leq f_D$, for every signed $t$ and model $h$,
\[
 H_{t,h}(u)\leq E(P_\theta(f_D+th))=\mathscr E(f_D+th),
 \qquad H_{0,h}(u)=\mathscr E(f_D).
\]
The two functions touch at zero. Divide by $t>0$ and $t<0$
separately, then differentiate using (13) and the finite derivative
(6). The opposite inequalities force $M(u)(h)=\mu_D(h)$.
Uniform density and equal mass identify the measures.
No ordered-measure continuity or attained inverse is used.
\end{proof}

\section{General weak-and-energy continuity}
\begin{theorem}\label{continuity}
If genuine $u_i\to u$ weakly and $E(u_i)\to E(u)>-\infty$,
then $M(u_i)\to M(u)$ weakly on the original compact carrier.
In particular this holds for every decreasing genuine finite-energy
net with finite-energy endpoint.
\end{theorem}
\begin{proof}
Normalize by generic values, which converge weakly; the constant
law preserves energy convergence and the outputs.
Their normalized costs are eventually bounded by some $B$.
Let $\lambda$ be any cluster of their positive mass-$V$ outputs.
At every fixed small signed $t$, the uniform genuine Taylor
estimate and the weak usc of (11) give
\[
 E(u)+t\lambda(h)-K_{B,c,h}t^2-|t|a_{B,h}/c
       \leq H_{t,h}(u).
\]
Divide separately by the two signs of $t$ and use (13).
First $t\to0$ at fixed $c$, then $c\to\infty$, gives
$\lambda(h)=M(u)(h)$. Density identifies every cluster, hence
the whole net. Decreasing genuine nets have coordinate convergence
and energy convergence by Lemma~\ref{outer}.
\end{proof}

\section{Algebraic nullness without relative integrability}
\begin{theorem}\label{null}
Every $M(u)$ gives every proper algebraic closed subset zero mass.
In particular $M(u)\{R=-\infty\}=0$.
\end{theorem}
\begin{proof}
Normalize $u$ with cost $B$, and take its nonpositive model
majorants. For a nonzero section of an actual effective divisor $S$
on $X$, put $v(S)=-\log|s_S|\geq0$, with value $+\infty$ on
$S^{\mathrm{an}}$. On the model of $D$, its total pullback consists
of an effective horizontal part and vertical coefficients
$b_E x_E(S)$. The finite positive deficient measure satisfies
\[
 \int v(S)\,d\mu_{D,c}
 \leq\langle\alpha_D(c)^n\rangle\cdot\rho_X^*S,
 \quad \alpha_D(c)=\rho_X^*L+D+cF.                      \tag{15}
\]
This uses nonnegativity of the volume derivative in the effective
horizontal direction. It is not a claim of vanishing intersection
with a generic fixed factor.

Fix $c>B/V$, and a fixed determination carrying $S$ and an ample
bundle $T$ with $aT-\rho_X^*S$ nef. Pull both tests to every
refinement. Because $D\leq0$, positive-product order against the
nef $T$ bounds the right side of (15) by
\[
 C_c=a\langle(\rho_X^*L+cF)^n\rangle\cdot T<\infty,
                                                        \tag{16}
\]
independently of $D$. The deficient result in the finite input gives
$0\leq\mu_D-\mu_{D,c}$ with mass at most $B/c$.
The original continuous function $b_N=\min(v(S),N)$ is
$N$ on $S^{\mathrm{an}}$. Consequently
\[
 \int b_N\,d\mu_D\leq C_c+NB/c.
\]
Pass the model-majorant net to $M(u)$ at this fixed continuous
test, divide by $N$, then $N\to\infty$. It follows that
$M(u)(S^{\mathrm{an}})\leq B/c$.
Now $c\to\infty$ proves nullness. Each proper closed subset of
projective $X$ is contained in an effective divisor, so the first
assertion follows. Finally choose an actual Kodaira identity
$L=aA+E_0$, with $A$ ample and $E_0$ effective.
Its section tensor with a generating system is a genuine minorant
$q_{E_0}\leq R$, so $\{R=-\infty\}\subset E_0^{\mathrm{an}}$.
\end{proof}

\section{Exact scope and remaining original obligations}
The preceding construction proves actual whole outer-energy and
outer-derivative statements from independently checked scalar and
finite inputs. No conditional S2 or M1 package has been inserted
as a premise. The touching proof makes model-envelope compatibility
independent of the later continuity proof.

We have not asserted that this outer derivative already equals
every unbounded physical moving-product definition. The quadratic
relative contraction, the full finite-list moving signatures,
all-test integrability, every negligible-set nullness, positive
genuine mixed identities and the intrinsic inverse remain separate
proof and review obligations. Algebraic nullness alone supplies
none of those conclusions. In particular the reference defect
$R-q_{T_m}$ in physical weighted limits is retained.
Original P16.5 remains open and no acceptance count changes.

The whole source of the new construction is written here. The
complete finite and scalar bodies \cite{Finite,Scalar} and their independent
round 161 scope binding supply precisely the two stated inputs.
No new primary-source retrieval, whole-primary reading, recursive
background reconstruction, or independent acceptance of this paper
is claimed. The existing user editor remains unchanged.
\begin{thebibliography}{9}
\bibitem{Finite} Research paper, round 159,
\emph{Absolute model estimates and a conditional attainment mechanism},
6 October 2026, Section 4.
\href{https://github.com/MingchenXia/nonarchimedean-pluripotential-theory/blob/2ec4386cbeec687e2a75a26b457ba9e7fd574aa0/research-checkpoints/continuation169/root-checks/public-P16.5-round159-partial-conditional-attainment-reader/paper.pdf}{Frozen full paper and source in GitHub checkpoint 169}.
\bibitem{Scalar} Research paper, round 160,
\emph{Single-net regularization and exact negligible-set transfer in big classes},
6 October 2026, Sections 2--4.
\href{https://github.com/MingchenXia/nonarchimedean-pluripotential-theory/blob/2ec4386cbeec687e2a75a26b457ba9e7fd574aa0/research-checkpoints/continuation169/root-checks/public-P16.5-round160-partial-negligible-transfer-reader/paper.pdf}{Frozen full paper and source in GitHub checkpoint 169}.
\bibitem{BFJ} S. Boucksom, C. Favre and M. Jonsson,
\emph{Differentiability of volumes of divisors and a problem of Teissier},
\href{https://arxiv.org/html/math/0608260v2}{arXiv:math/0608260v2}.
\bibitem{BJ} S. Boucksom and M. Jonsson,
\emph{Global pluripotential theory over a trivially valued field},
\href{https://arxiv.org/pdf/1801.08229v4}{arXiv:1801.08229v4}.
\end{thebibliography}
\end{document}
